
LiaScript is a Markdown-based language for creating interactive Open Educational Resources that runs entirely client-side, enabling courses to execute in the learner’s browser without server infrastructure. This sovereignty-by-design approach, due to eliminating accounts, server dependencies, and data-collection, gives educators full control over their content. However, this also eliminates the telemetry that would grant developers more insight into how pedagogical features are used in practice.

This paper addresses that gap through a corpus analysis of 3,317 validated LiaScript courses from public GitHub repositories (March 2026), contributed by 314 authors; of these, 2,973 courses with a clear education-level classification enter our group-level comparison. We introduce a didactically motivated taxonomy spanning presentation, interaction, reuse, and embedding features, and compare adoption across different authoring contexts --- core developers, a school-focused initiative, and independent community authors in higher education and school settings. Results show that 89\,\% of the designed feature space is adopted in measurable numbers, yet each community explores only its own subset: the challenge is not that features go unused, but that usage patterns vary strongly across contexts. Authoring infrastructure shapes adoption as strongly as pedagogical needs. We identify three adoption gap classes --- documentation deficits, workflow barriers, and configuration lock-in --- and derive context-sensitive implications for future development.
The study also demonstrates a replicable corpus methodology for understanding usage in OER that, by design, collects no user data.